\documentclass[11pt, a4paper, UTF8]{chinese-erj}

% --- 必须的宏包 ---
\usepackage{amsmath, amssymb, mathtools}

% --- 动态引入本篇稿件的参考文献库 ---
\addbibresource{AERub_lemon.bib}

% --- 动态定义出版信息 ---
\header{吉米奈、天马行空：逆向选择与“柠檬”相亲市场}
\pubinfo{2026年第2卷第1期} 

% --- 自定义无标号底部脚注，用于输出“收稿日期、作者简介”等 ---
\makeatletter
\newcommand{\bottominfo}[1]{%
  \begingroup
  % 重定义脚注格式：取消前面的上标占位和缩进，全左对齐，字号设为 small
  \def\@makefntext##1{\noindent\small ##1}% 
  \renewcommand\thefootnote{}%
  \footnotetext{#1}%
  \endgroup
}
\makeatother

% ==========================================
% 核心修改：使用模板原生的标题、作者与致谢命令
% ==========================================
\title{逆向选择与“柠檬”相亲市场\\——中国特色婚恋市场的阿克洛夫模型\thanks{吉米奈，谷歌大中华区人工智能研发中心，邮政编码：100000，电子信箱：gemini@google.com；天马行空（通讯作者），小红书大学经济学院，邮政编码：200000，电子信箱：skyhorse@rednote.edu。本文受自燃科学鸡精委重点项目（项目批准号：72636007）“人类迷惑行为的经济学解释”资助。作者感谢匿名审稿专家的宝贵建议，当然文责自负。}}
\author{吉米奈 \and 天马行空}
\date{} % 隐藏默认日期

\begin{document}

% 这一句会自动生成标题，并将 \thanks 里的内容变成第一页底部的脚注！
\maketitle

% ==========================================
% 核心修改：使用模板原生的摘要与关键词环境
% ==========================================
\begin{abstract}
本文将经典的信息不对称与“柠檬市场”模型引入中国特色相亲市场。研究发现，由于相亲市场存在严重的信息不对称（如高度美颜的照片、隐瞒的性格缺陷及夸大的经济实力），相亲对象的真实质量对于匹配方而言是不可观测的。在这种机制下，高质量单身青年（“好车”）会因为无法获得与之匹配的估值而主动退出相亲市场；而带有隐性瑕疵的单身青年（“柠檬/破车”）则会大量涌入。数学推导表明，随着多轮博弈的进行，相亲市场的平均质量将持续恶化，最终陷入“劣币驱逐良币”的死循环，导致整个相亲市场彻底沦为“奇葩收容所”。本文为当代青年“越相亲越绝望”的普遍社会现象提供了严谨的微观经济学证明。

\keywords{柠檬市场 \quad 逆向选择 \quad 信息不对称 \quad 相亲经济学 \quad 阿克洛夫模型}
\end{abstract}

% --- 正文 ---
\section{引言}

在传统的中国式家长眼中，相亲市场是一个高效的资源配置平台（类似完美的纳斯达克交易中心），能够迅速撮合男女双方实现效用最大化。然而，当代青年的实际相亲体验却呈现出一种令人绝望的“开盲盒”特征：长辈口中“老实本分、工作稳定”的优质资产，在实际见面后往往被证实为情商堪忧、充满隐性雷区的“不良资产”。

为什么相亲局上遇到“正常人”的概率越来越低？本文拒绝使用心理学或社会学的感性分析，而是直接引入 \textcite{akerlof1970} 的经典“柠檬市场（The Market for Lemons）”模型，通过数学推导证明：相亲市场沦为“奇葩聚集地”不仅不是偶然，而是信息不对称下的必然市场均衡。

\section{信息不对称与“伪装成本”}

相亲市场的核心痛点在于极端的信息不对称（Information Asymmetry）。

在自由恋爱市场（Primary Market）中，双方通过长期的同学、同事关系或共同社交圈，能够以极低的成本获取对方的真实质量信号（如性格、三观、真实素颜等）。然而，相亲市场是一个典型的“二手车交易市场”（Secondary Market）。推荐人（媒人或亲戚）往往为了促成交易而粉饰太平；同时，随着“美颜滤镜”技术的普及和朋友圈人设的工业化包装，隐藏自身缺陷的“伪装成本”趋近于零。

作为买方（参与相亲的一方），你只能看到对方的“预期质量”，而无法在见面（甚至结婚）前观测到其“真实质量”。

\section{逆向选择的数学推导}

假设相亲市场中潜在参与者的真实质量（包括性格、外貌、经济实力综合得分）为 $q$。为简化模型，假设 $q$ 在区间 $[0, 1]$ 上均匀分布，即 $q \sim U[0, 1]$。

买方（相亲对象的审视者）无法观测到具体的 $q$，只能根据市场平均水平给出自己的最高出价/投入意愿 $P$（如愿意付出的彩礼、嫁妆、情绪价值及时间成本）。在风险中性假设下，买方的出价等于市场上所有愿意参与相亲的人的预期质量：
\begin{equation}
    P = E[q]
\end{equation}

卖方（相亲者本人）清楚自己的真实质量 $q$。只有当市场给出的估值（对方的投入/条件）大于或等于自己的真实质量时，卖方才愿意留在相亲市场中。即，只有当 $q \le P$ 时，参与者才会去相亲。

这就触发了致命的逆向选择（Adverse Selection）机制：由于只有质量 $q \le P$ 的人才会留在市场里，买方所面对的“预期质量”将发生断崖式下跌：
\begin{equation}
    E[q \mid q \le P] = \frac{P}{2}
\end{equation}

由于买方也是理性的，他们会立刻意识到这一轮市场质量的下降，并将自己的出价 $P$ 调整为 $\frac{P}{2}$。这又会导致质量在 $[\frac{P}{2}, P]$ 区间的相对优质青年觉得“对方条件太差，配不上自己”，从而愤怒地退出相亲市场。

\section{均衡结果：市场的全面崩盘}

上述博弈将进入一个死循环：
\begin{enumerate}
    \item 市场给出平均估价 $P$。
    \item 质量高于 $P$ 的“优质单身青年”（好车）觉得受到了侮辱，退出市场。
    \item 市场平均质量暴跌至 $\frac{P}{2}$。
    \item 估价随之暴跌至 $\frac{P}{2}$。
    \item 质量高于 $\frac{P}{2}$ 的中等青年也退出市场\dots
\end{enumerate}

极限状态下，市场的出价和质量将收敛于：
\begin{equation}
    \lim_{n \to \infty} P_n = 0
\end{equation}

\textbf{定理1（相亲市场的阿克洛夫崩盘）：} 在缺乏有效信号传递机制的相亲市场中，由于高昂的信息不对称，高质量个体（正常人）将不可避免地被全部挤出市场。最终，相亲市场的主体将仅仅由 $q \to 0$ 的极端劣质资产（奇葩/柠檬）构成。交易量萎缩，市场发生绝对失灵。

\section{结论与政策建议}

本文用冷酷的数学公式证明了当代年轻人的一句口头禅：“相亲市场里没有正常人，因为正常人早就在自由恋爱市场里内部消化了。”

面对这一帕累托极度无效率的市场，我们提出以下经济学建议：第一，引入强信号传递机制（Signaling）\parencite{spence1973}。相亲时建议首次见面直接出示“中国人民银行个人征信报告”、“近三年三甲医院全身体检报告”以及“无美颜身份证原件”，这些高成本强信号有助于打破信息不对称；第二，拥抱“主动退市”策略。在一个注定充满“柠檬”的市场里，尽早认清现实，拒绝相亲安排，将精力投入自身人力资本积累，是唯一能够实现自身效用最大化的纳什均衡。

% --- 输出参考文献 ---
\erjref
\printbibliography[heading=none]

\end{document}